\documentclass[11pt,letterpaper]{article}
\usepackage[T1]{fontenc}
\usepackage[utf8]{inputenc}
\usepackage{lmodern}
\usepackage[margin=1in,headheight=14pt]{geometry}
\usepackage{amsmath,amssymb,amsthm,mathtools}
\usepackage{booktabs,tabularx,array}
\usepackage{enumitem}
\usepackage{microtype}
\usepackage{xcolor}
\usepackage{fancyhdr}
\usepackage{hyperref}
\definecolor{inkblue}{RGB}{25,60,95}
\hypersetup{colorlinks=true,linkcolor=inkblue,citecolor=inkblue,urlcolor=inkblue,
  pdftitle={ALCI RCC5: Formal Audit and Partial Results},
  pdfauthor={OpenAI Codex; prepared for Michael Wessel}}
\pagestyle{fancy}
\fancyhf{}
\fancyhead[L]{\small\textsc{ALCI\_RCC5: audit and partial results}}
\fancyhead[R]{\small 7 October 2026}
\fancyfoot[C]{\thepage}
\renewcommand{\headrulewidth}{0.3pt}
\setlength{\parindent}{0pt}
\setlength{\parskip}{5pt plus 1pt}
\setlist{itemsep=2pt,topsep=4pt}
\setcounter{tocdepth}{1}
\newcommand{\logic}{\ensuremath{\mathcal{ALCI}_{\mathrm{RCC5}}}}
\newcommand{\EQ}{\mathsf{EQ}}
\newcommand{\PP}{\mathsf{PP}}
\newcommand{\PPI}{\mathsf{PPI}}
\newcommand{\PO}{\mathsf{PO}}
\newcommand{\DR}{\mathsf{DR}}
\newcommand{\Rel}{\mathcal R}
\newcommand{\Sat}{\operatorname{Sat}}
\newcommand{\er}{\operatorname{erase}}
\newcommand{\cl}{\operatorname{cl}}
\newcommand{\cone}{\operatorname{cone}}
\newcommand{\Bod}{\operatorname{Bod}}
\newcommand{\code}[1]{\texttt{#1}}
\newcommand{\inc}{\mathrel{\|}}
\newtheorem{theorem}{Theorem}[section]
\newtheorem{lemma}[theorem]{Lemma}
\newtheorem{proposition}[theorem]{Proposition}
\newtheorem{corollary}[theorem]{Corollary}
\theoremstyle{definition}
\newtheorem{definition}[theorem]{Definition}
\newtheorem{example}[theorem]{Example}
\theoremstyle{remark}
\newtheorem{remark}[theorem]{Remark}
\newcolumntype{Y}{>{\raggedright\arraybackslash}X}

\title{\vspace{-1.2em}\textbf{\logic: Audit of the Decidability Claim\\[3pt]
and Two Partial Results Beyond the\\[3pt]
\texorpdfstring{$\forall\PO$}{forall-PO}-Free Fragment}}
\author{Technical report prepared for Michael Wessel\\
\small OpenAI Codex}
\date{7 October 2026}

\begin{document}
\maketitle
\thispagestyle{empty}
\begin{abstract}
This report examines the current project snapshot supplied as
\texttt{alcircc5-master(1).zip}, using its main Markdown pages and their
current paper and Lean references. No gap was found in the present
$\forall\PO$-free concept-satisfiability argument or in the inspected
dependencies of its Lean capstones. This was a source audit, not a fresh Lean
kernel run. Unrestricted \logic\ remains unresolved.

Two mathematical results are proved. First, for a fixed strict part order,
disjointness has a unique largest admissible interpretation: two points are
disjoint exactly when they have no common lower bound. Principal down-sets
realize it, and the transformation preserves every NNF concept without
$\exists\PO$ or $\forall\DR$. Second, for NNF concepts without
$\exists\PP$, $\exists\PO$, or $\forall\DR$, replacing each
$\forall\PO$ restriction by $\top$ preserves satisfiability. A fresh
downward witness forest proves the second result and yields a decidable
syntactic extension of the certified fragment. Explicit counterexamples
delimit both constructions. These new proofs have not been formalized in Lean
and do not decide the unrestricted logic.
The report concludes with a separately identified, subjective outlook on
the prospect of an AI-assisted resolution over the next two to five years.
\end{abstract}

\section{Scope and conclusions}
\label{sec:scope}

The authoritative status documents for this review are the supplied
\path{README.md}, the approach pages for the cone scheme and Lean/F6, and
the current sources \path{papers/pofree_fragment_arxiv.tex} and
\path{papers/overview_arxiv.tex}. The relevant formal development is
\path{formal/POFreeLift.lean}. Historical failures were used only to
identify a potential source of error; they were not treated as current proofs.

The subject is \emph{single-concept satisfiability} under the abstract RCC5
composition-table semantics with strong equality. The roles are the five
atomic RCC5 relations. The report does not assert a decision result for
general TBoxes or for a fixed domain containing \emph{all} regions of a
particular space. All syntactic restrictions below are imposed after
conversion to negation normal form (NNF).

\begin{center}
\small
\begin{tabularx}{\linewidth}{@{}p{0.37\linewidth}Y@{}}
\toprule
Claim & Assessment in this review\\
\midrule
$\forall\PO$-free satisfiability & No gap found in the current proof or
inspected Lean dependencies. No fresh kernel build was performed.\\
Full unrestricted satisfiability & Still open. The certificate-completeness
premise required by the full-language reduction is unproved.\\
Principal-down-set normalization & Proved here for the NNF fragment
without $\exists\PO$ and $\forall\DR$. No decidability conclusion for
that whole fragment.\\
PO erasure & Proved here when $\exists\PP$, $\exists\PO$, and
$\forall\DR$ are absent. Gives a decidable syntactic extension.\\
\bottomrule
\end{tabularx}
\end{center}

The two new results are relative to the established RCC5 semantics. They
are mathematical results of this review, not claims of literature priority.
They have been independently checked for the relation orientations and
boundary cases, but have not been added to the Lean artifact.

\section{Semantics and the order presentation}
\label{sec:semantics}

Let $\Rel=\{\EQ,\PP,\PPI,\PO,\DR\}$. An RCC5 frame has a nonempty
carrier $W$ and a total relation labelling $\rho:W^2\to\Rel$, with
identity as $\EQ$, the usual converse operation, and the RCC5 composition
condition. Here $x\,\PP\,y$ means that $x$ is a \emph{proper part of}
$y$; thus a $\PP$ witness is a superpart and a $\PPI$ witness is a part.

An interpretation adds a valuation $V(A)\subseteq W$ for each concept name.
NNF concepts have grammar
\[
C ::= \top\mid\bot\mid A\mid\neg A\mid C\sqcap C\mid C\sqcup C
       \mid\exists R.C\mid\forall R.C,\qquad R\in\Rel.
\]
The modal clauses are
\begin{align*}
\mathcal M,x\models\exists R.D
 &\iff \exists y\in W\,[\rho(x,y)=R\ \land\ \mathcal M,y\models D],\\
\mathcal M,x\models\forall R.D
 &\iff \forall y\in W\,[\rho(x,y)=R\ \Longrightarrow\ \mathcal M,y\models D].
\end{align*}
Satisfiability means truth at a point of some such interpretation.

\begin{definition}[Ordered-disjoint structure]
An ordered-disjoint structure is $(W,<,\#)$ where $<$ is a strict partial
order and $\#$ is symmetric and irreflexive, with the following properties.
Writing $\leq$ for the reflexive closure of $<$,
\begin{align}
x<y&\ \Longrightarrow\ \neg(x\#y),\label{eq:separate}\\
x\#y,\quad u\leq x,\quad v\leq y
 &\ \Longrightarrow\ u\#v.\label{eq:downward}
\end{align}
Its induced labelling gives $\EQ$ to equality, $\PP$ to $<$,
$\PPI$ to $>$, $\DR$ to $\#$, and $\PO$ to the remaining pairs.
\end{definition}

\begin{theorem}[Established RCC5 normal form]
\label{thm:normal}
The induced labelling of every ordered-disjoint structure is an RCC5
frame. Conversely, reading $<$ as $\PP$ and $\#$ as $\DR$ in any
strong-EQ RCC5 frame yields an ordered-disjoint structure with the original
labelling. The abstract frames and families of nonempty, pairwise distinct
sets have the same concept satisfiability.
\end{theorem}

This theorem is an established input, proved in the current project
\cite{fragment,lean}. Transitivity comes from
$\PP\circ\PP=\{\PP\}$; downward inheritance follows from
$\PP\circ\DR=\{\DR\}$ and $\DR\circ\PPI=\{\DR\}$.
The converse is the checked composition-table case analysis. The project's
general set representation retains overlap even when no common part is
represented in the carrier; this distinction will matter below.

Under set semantics, the domain is a chosen family of nonempty sets.
Equality, strict inclusion, converse inclusion and disjointness interpret
the first four applicable cases; $\PO$ means nonempty intersection and
incomparability under inclusion. A shared point of the ambient set universe
need not itself occur as a named region in the domain.

\section{Audit of the existing decidability argument}
\label{sec:audit}

\subsection{Finite signatures and unconditional completeness}

For a fixed NNF concept $C_0$, the current cone procedure takes support
types $T\subseteq\cl(C_0)$ and signatures $q=(T,S)$, where $S$ is a set
of candidate types allowed strictly below $q$. Let
$\cone(q)=\{T\}\cup S$ and
$\Bod_R(T)=\{D:\forall R.D\in T\}$.
Static admissibility checks the local Boolean and EQ conditions, and
for each $U\in S$ requires
\[
\Bod_{\PP}(U)\subseteq T,
\qquad
\Bod_{\PPI}(T)\subseteq U.
\]
A target $q'=(T',S')$ for a non-EQ demand $\exists R.D\in T$ must
contain $D$ in $T'$ and satisfy the appropriate compatibility condition:
\begin{center}
\small
\begin{tabularx}{\linewidth}{@{}lY@{}}
\toprule
$R$ & Compatibility condition\\
\midrule
$\PP$ & $\cone(q)\subseteq S'$.\\
$\PPI$ & $\cone(q')\subseteq S$.\\
$\DR$ & For all $U\in\cone(q)$ and $U'\in\cone(q')$,
$\Bod_{\DR}(U)\subseteq U'$ and $\Bod_{\DR}(U')\subseteq U$.\\
$\PO$ & $\Bod_{\PO}(T)\subseteq T'$ and
$\Bod_{\PO}(T')\subseteq T$.\\
\bottomrule
\end{tabularx}
\end{center}
Pruning removes signatures that lack a compatible target for some demand.
The operator is monotone and reductive on a finite initial signature set.

In any actual model, take the complete closure type of a point and the set
of closure types realized at its proper parts. These signatures are
statically admissible and support one another using the model's witnesses.
They consequently survive every pruning round. This direction does
\emph{not} assume $\forall\PO$-freeness. The checked declaration is
\code{coneScheme\_complete}.

\subsection{Fresh occurrences and fragment soundness}

The converse unfolds the surviving finite control states to fresh
occurrences. Part births generate a strict order; DR births generate
disjointness together with all its downward consequences. Residual pairs
are PO. The height and vertical-component invariants establish a genuine
ordered-disjoint structure. Cone propagation handles all PP and PPI
universals, and the cross-cone test at every DR birth handles all DR
universals. Witness children serve existential demands.

The truth lemma uses the fragment hypothesis precisely when it excludes
$\forall\PO.D$ from occurrence labels. Most PO pairs were not created
as witness edges: they are residual pairs between branches. A universal
over PO would constrain these pairs too. This step appears in
\code{unf\_truth}, and gives \code{coneScheme\_sound}. The finite
fixed point then yields \code{decidableSat\_cone}.

The audit found no missing premise in this current argument. In
particular, the DR test compares entire downward cones, and the unfolding
does not equate two occurrences simply because their signatures agree.
Both details are necessary. The source also proves normalization from raw
negation in both polarities. The raw fragment excludes positive
$\forall\PO$ and negative $\exists\PO$; merely checking the spelling
of modalities before NNF would be incorrect.

\subsection{Verification strength and limits}

The audit inspected the capstone statements and their proof dependencies,
the semantics, the raw polarity test, and the abstract-to-set equivalence.
No executable \code{sorry}, \code{admit}, or new user axiom was identified
in the current formal files. A separate finite enumeration recovered all
25 displayed RCC5 composition cells from nonempty subsets of a
four-element universe.

Lean was unavailable in the review environment, and the archive contained
no compiled \code{.olean} artifacts. Accordingly, this report does not
claim an independent kernel recompilation or fresh axiom output. The
repository reports such certification; the present review found no
source-level defect in the relevant argument. The theorem
\code{satisfiable\_iff\_set} concerns arbitrary set families. The
paper's further regular-closed-region corollary uses the representation
theorems of Lutz and Wolter \cite{lutzwolter} externally.

\section{Maximal disjointness and principal down-sets}
\label{sec:downsets}

Fix a nonempty strict partial order $(W,<)$. Write $x\inc y$ when $x,y$
are incomparable under its reflexive closure. Define
\begin{equation}
x\mathrel{\#_{\max}}y
\quad\Longleftrightarrow\quad
\neg\exists z\in W\ (z\leq x\ \land\ z\leq y).
\label{eq:maxdr}
\end{equation}

\begin{lemma}[Largest admissible disjointness]
\label{lem:maxdr}
The relation $\#_{\max}$ is the unique largest relation, under inclusion,
which makes $(W,<,\#)$ an ordered-disjoint structure.
\end{lemma}
\begin{proof}
Symmetry is immediate. Irreflexivity holds because $x$ is a lower bound
of $x,x$. If $x\leq y$, then $x$ is a common lower bound, so comparable
pairs are not in $\#_{\max}$. For downward inheritance, suppose
$x\#_{\max}y$, $u\leq x$ and $v\leq y$. A common lower bound of $u,v$
would, by transitivity, be a common lower bound of $x,y$, a contradiction.
Thus $\#_{\max}$ is admissible.

If $\#$ is any admissible relation and $x\#y$, a common lower bound
$z\leq x,y$ would yield $z\#z$ by downward inheritance, contradicting
irreflexivity. Hence $\#\subseteq\#_{\max}$, proving maximality and
uniqueness of the largest relation.
\end{proof}

For each $x\in W$, put $P_x=\{z\in W:z\leq x\}$.

\begin{lemma}[Exact set representation]
\label{lem:principal}
The family $(P_x)_{x\in W}$ consists of nonempty, pairwise distinct sets,
and
\[
P_x\subseteq P_y\iff x\leq y,
\qquad
P_x\cap P_y=\varnothing\iff x\#_{\max}y.
\]
Consequently its RCC5 labelling preserves the order and gives
\begin{equation}
\PO_{\min}(x,y)
\quad\Longleftrightarrow\quad
x\inc y\ \land\ \exists z\in W\ (z\leq x\land z\leq y).
\label{eq:minpo}
\end{equation}
\end{lemma}
\begin{proof}
The element $x$ belongs to $P_x$. If $x\leq y$, transitivity gives
$P_x\subseteq P_y$; conversely, that inclusion puts $x$ in $P_y$.
Equality of the two sets therefore implies $x\leq y$ and $y\leq x$,
hence $x=y$. Their intersection consists exactly of the common lower
bounds. The RCC5 set clauses now give~\eqref{eq:minpo}.
\end{proof}

For an RCC5 interpretation $\mathcal M$, let $\mathcal M_{\max}$ be its
principal-down-set interpretation, identified with the same index set $W$,
and with the same unary valuations. Its order is unchanged and its
disjointness is $\#_{\max}$.

\begin{theorem}[Truth preservation under maximal disjointness]
\label{thm:normalization}
Let $C$ be an NNF concept with no occurrence of $\exists\PO$ or
$\forall\DR$, including inside modal bodies. For every RCC5
interpretation $\mathcal M$ and $x\in W$,
\[
\mathcal M,x\models C\quad\Longrightarrow\quad
\mathcal M_{\max},x\models C.
\]
Therefore satisfiability for this fragment can be restricted to
principal-down-set interpretations.
\end{theorem}
\begin{proof}
Lemmas~\ref{lem:maxdr}--\ref{lem:principal} show that EQ, PP and PPI are
unchanged, old DR edges remain DR, and only old PO edges can change, each
either remaining PO or becoming DR. Induct on $C$. Literals, constants and
Boolean connectives follow from the unchanged valuation and the induction
hypothesis. For EQ, PP and PPI modalities, use their unchanged successor
sets and the induction hypothesis on the body. An old DR existential
witness is still a DR witness and still satisfies its body. Every new PO
successor was an old PO successor, so a PO universal continues to hold by
the induction hypothesis. The two modal cases that could fail are
precisely the two excluded constructors. The normalized interpretation is
itself an RCC5 interpretation, giving the converse implication for
\emph{existence} of a satisfying model without any additional argument.
\end{proof}

\begin{remark}[Normalization is not yet a decision procedure]
The order can still force PO. For example,
\[
\exists\PP.(A\sqcap\forall\PP.A)
\ \sqcap\
\exists\PP.(\neg A\sqcap\forall\PP.\neg A)
\]
forces distinct incomparable superparts $a,b$ of the root: either order
between them violates one of the PP universals. Their shared proper part
excludes DR, so they must be PO. The normal form retains such overlaps.
\end{remark}

\section{A decidable fragment by erasing PO universals}
\label{sec:erasure}

\begin{definition}[The downward-witness fragment]
Let $\mathcal D$ be the class of NNF concepts with no occurrences of
$\exists\PP$, $\exists\PO$ or $\forall\DR$. All other constructors
are permitted, including nested $\forall\PO$.
Define $\er$ recursively by
\[
\er(\forall\PO.D)=\top,
\]
and homomorphically on every other constructor. In particular, an erased
PO restriction loses its entire body. Then $\er(C)$ is $\forall\PO$-free
and $|\er(C)|\leq |C|$.
\end{definition}

\begin{lemma}[PO-empty witness-forest property]
\label{lem:forest}
Every satisfiable NNF concept $E$ without $\forall\PO$, $\exists\PP$,
$\exists\PO$, or $\forall\DR$ has a countable model in which PO is empty.
\end{lemma}
\begin{proof}
Let $\mathcal M,x_0\models E$. We construct occurrences in finitely
many stages per depth. Each occurrence $u$ carries a source point
$f(u)\in W^{\mathcal M}$. Begin with $u_0$ and $f(u_0)=x_0$.
For each already generated occurrence $u$ and every existential
subformula occurrence $e=\exists R.D$ of $E$ with
$\mathcal M,f(u)\models e$ and $R\neq\EQ$, choose a witness $y$ in
$\mathcal M$ and generate a fresh occurrence $u\cdot e$ with
$f(u\cdot e)=y$. There are finitely many such requirements per occurrence,
so the set $U$ of all generated finite histories is countable. The only
non-EQ roles needing witnesses are PPI and DR.

Retain the parent edge between $u$ and $u\cdot e$ only when $e$ has role
PPI, and put $u\cdot e<u$. A DR child starts a new component with no
vertical edge to its generating occurrence. The transitive closure of
these retained edges is a strict order. Its components are downward
rooted trees: every non-root has exactly one immediate superpart, and
every point has a finite chain of ancestors to its component root.

Set $u\#v$ exactly when $u,v$ are incomparable. This is downward
hereditary. Indeed, if $u',v'$ lie below incomparable $u,v$ but are
comparable, one of $u',v'$ lies below both $u,v$. Its ancestors form a
chain, forcing $u,v$ to be comparable, a contradiction. Thus the order
and $\#$ form an ordered-disjoint structure, with no residual PO pairs.
Equivalently, the principal down-sets of its nodes are nested along
branches and disjoint between incomparable nodes, directly supplying a
set representation. Let $\mathcal N$ be this frame with
$u\in V^{\mathcal N}(A)$ iff $f(u)\in V^{\mathcal M}(A)$.

Whenever $u<v$ in $\mathcal N$, the source points satisfy
$f(u)\,\PP\,f(v)$ in $\mathcal M$: each edge has this property by
the selected PPI witness, and the property composes transitively. Hence
the map $f$ preserves every new PP and PPI edge in its original direction.
It also preserves equality. A generated DR child lies in a fresh
component and therefore is DR from its generating point.

We now prove simultaneously, for every subformula $D$ of $E$ and every
occurrence $u$, the implication
\begin{equation}
\mathcal M,f(u)\models D\quad\Longrightarrow\quad
\mathcal N,u\models D.
\label{eq:truthtransfer}
\end{equation}
Induct on $D$. Literals and Boolean cases follow from the pulled-back
valuation and the induction hypothesis. Non-EQ existential cases are
served by their generated children, whose sources satisfy the body.
EQ modalities reduce to the body at the same point. For a PP or PPI
universal, each new successor has the same relation between its source
points, so the old universal and the induction hypothesis apply. No other
universal or existential case remains under the stated restrictions.
Applying~\eqref{eq:truthtransfer} to $E$ at $u_0$ proves the result.
\end{proof}

\begin{theorem}[PO-erasure equivalence]
\label{thm:erasure}
For every $C\in\mathcal D$,
\[
\Sat(C)\quad\Longleftrightarrow\quad\Sat(\er(C)).
\]
\end{theorem}
\begin{proof}
In NNF all compound constructors are monotone in their concept arguments.
Replacing an entire $\forall\PO$ subformula by $\top$ therefore
preserves satisfaction, proving the forward implication.

Conversely, $\er(C)$ satisfies the hypotheses of Lemma~\ref{lem:forest}.
If it is satisfiable, it has a model $\mathcal N$ with empty PO. In any
such model, for every concept $D$ and every point $u$,
\[
\mathcal N,u\models D\quad\Longleftrightarrow\quad
\mathcal N,u\models\er(D).
\]
This follows by structural induction: a PO universal and its replacement
$\top$ are both true, and all remaining constructors commute with
erasure. In particular, the point satisfying $\er(C)$ satisfies $C$.
\end{proof}

\begin{corollary}[Effective extension of the existing decision result]
\label{cor:decision}
Satisfiability for $\mathcal D$ is decidable. So is satisfiability for
the union of $\mathcal D$ with the published $\forall\PO$-free fragment.
This union is a strict syntactic extension of that fragment.
\end{corollary}
\begin{proof}
Both syntactic membership tests are effective. For $C\in\mathcal D$,
compute $\er(C)$ and apply the established $\forall\PO$-free decision
procedure. Use the established procedure directly on the original
fragment. The formula $\forall\PO.A$ belongs to $\mathcal D$ but is
not $\forall\PO$-free.
\end{proof}

This extension does not add an irreducible PO-universal satisfiability
constraint: those restrictions are erasable within $\mathcal D$. It is
nevertheless a decision theorem about inputs outside the original
syntactic fragment. Its new model-construction proof is not yet Lean
checked. The same doubly exponential upper bound stated for the literal
cone procedure applies after erasure, since erasure does not increase size;
no new lower bound is asserted.

\section{Counterexamples delimiting the constructions}
\label{sec:boundaries}

Each excluded constructor in Theorem~\ref{thm:erasure} is necessary for
that unrestricted erasure rule. In the following examples, erasure is
satisfiable while the original concept is unsatisfiable.

\subsection{Permitting an existential PO restriction}

The immediate example is
\[
C_{\exists\PO}=\exists\PO.\top\sqcap\forall\PO.\bot.
\]
It is unsatisfiable, whereas its erasure $\exists\PO.\top$ is
satisfied by two overlapping, incomparable sets.

There is also a sharper boundary for Theorem~\ref{thm:normalization}:
\begin{equation}
C_{\mathrm{optional}}=
\forall\PPI.A\sqcap\exists\PO.(\forall\PPI.\neg A).
\label{eq:optional}
\end{equation}
The two-region domain $x=\{0,1\}$, $y=\{1,2\}$ satisfies this at $x$;
neither region has a represented proper part. In any satisfying model,
a common lower bound of $x$ and its PO witness $y$ would be a strict
part of both, forcing both $A$ and $\neg A$. Thus the overlap must
lack a represented common part. Such a pair is DR under principal
down-sets. Consequently~\eqref{eq:optional} has no principal-down-set
model at all.

\subsection{Permitting an existential PP restriction}

Consider
\begin{align}
C_{\exists\PP}={}&
\exists\PP.(A\sqcap\forall\PP.A\sqcap\forall\PO.A)
\notag\\[-2pt]
&\sqcap\exists\PP.(\neg A\sqcap\forall\PP.\neg A).
\label{eq:ppboundary}
\end{align}
At a putative root $x$, choose the witnesses $a,b$ with $A(a)$ and
$\neg A(b)$. They are distinct. If $a\,\PP\,b$, the universal at
$a$ forces $A(b)$; if $b\,\PP\,a$, the universal at $b$ forces
$\neg A(a)$. They cannot be DR because their shared proper part $x$
would then be disjoint from itself. Thus $a\,\PO\,b$, but the PO
universal at $a$ again forces $A(b)$, a contradiction.

Its erasure has the three-region model
\[
x=\{0\},\qquad a=\{0,1\},\qquad b=\{0,2\},\qquad V(A)=\{a\}.
\]
The two PP universals are vacuous at the maximal regions $a,b$. This
shows why upward witnesses can prevent a PO-empty forest construction.

\subsection{Permitting a universal DR restriction}

Consider
\begin{align}
C_{\forall\DR}={}&
\exists\PPI.(A\sqcap\forall\PPI.A\sqcap\forall\DR.A
                         \sqcap\forall\PO.A)
\notag\\[-2pt]
&\sqcap\exists\PPI.(\neg A\sqcap\forall\PPI.\neg A).
\label{eq:drboundary}
\end{align}
Its two proper-part witnesses $a,b$ are distinct. If $a<b$, the PPI
universal at $b$ contradicts $A(a)$; if $b<a$, the PPI universal at
$a$ contradicts $\neg A(b)$. A DR relation contradicts the DR
universal at $a$, and a PO relation contradicts its PO universal.
Every RCC5 value is excluded, so the concept is unsatisfiable.

The erasure has the model
\[
x=\{0,1,2\},\qquad a=\{0,1\},\qquad b=\{1,2\},\qquad V(A)=\{a\}.
\]
All displayed PPI and DR universals at $a,b$ are vacuous. Moreover,
\emph{every} model of this erased concept forces $a,b$ to be PO:
the preceding case split still excludes EQ, PP, PPI and DR. A common
lower bound would satisfy both $A$ and $\neg A$ by their PPI universals.
Thus the erased concept also has no principal-down-set model. This proves
that admitting arbitrary $\forall\DR$ breaks the normalization theorem
even if $\exists\PO$ remains absent.

\section{Consequences for the unrestricted problem}
\label{sec:full}

The new normalization distinguishes two sources of overlap. For fixed
order, an incomparable pair with a represented common lower bound is
necessarily PO; this overlap survives every admissible choice of DR.
Other PO pairs are optional relative to the order, although concepts
can require them, as~\eqref{eq:optional} and the erased form
of~\eqref{eq:drboundary} demonstrate. A complete construction for the
unrestricted language must accommodate both kinds and the universal
constraints on all resulting pairs.

The current cone scheme has no such complete joint control. Its
full-language completeness direction provides a sound UNSAT test, but
acceptance is not sufficient. The current focused paper records
\begin{align*}
C_{\mathrm{bad}}={}&\exists\PPI.(\exists\PP.A)
  \sqcap\forall\PP.\neg A\sqcap\forall\PPI.\neg A
  \sqcap\neg A\sqcap\forall\PO.\neg A,\\
C_{\mathrm{joint}}={}&
\exists\PP.(\forall\PPI.A\sqcap\forall\PO.A)
  \sqcap\exists\PO.\neg A.
\end{align*}
Both are unsatisfiable but accepted by that one-sided test. The first
excludes all four possible relations between the root and an $A$-point
sharing a part; the second excludes both possibilities in
$\PPI\circ\PO=\{\PPI,\PO\}$ between its two witnesses.

On the separate full-language certificate route, the current state is
\[
\begin{aligned}
\text{F6}
&\ \overset{?}{\Longrightarrow}\ \text{refined uniformization}\\
&\ \overset{?}{\Longrightarrow}\ \text{complete candidate lists}
\ \Longrightarrow\ \text{decidability}.
\end{aligned}
\]
Only the last implication, with the required premise supplied, is the
certified conditional decision result described by the current overview.
The two new theorems do not discharge either question-marked implication.

The next useful target is a finite representation theorem, or a
counterexample to a precisely defined representation scheme, that handles
forced common-part overlaps together with required optional overlaps.
A pair-matrix extension must coordinate the relations among witnesses;
checking its success on $C_{\mathrm{bad}}$ and $C_{\mathrm{joint}}$ is a
necessary diagnostic, not a completeness proof. Conversely, disproving
a particular width bound would not by itself establish undecidability;
that would require a reduction from an undecidable problem.

\textbf{Final assessment.} The present $\forall\PO$-free argument
survived this audit. The new results give a controlled order-only
normalization and a further decidable syntactic slice. No proof of
decidability or undecidability of unrestricted \logic\ was obtained.

\section{Outlook: could more capable AI resolve the problem soon?}
\label{sec:outlook}

\paragraph{Judgment and horizon.}
Taking ``near future'' to mean approximately the next two to five years,
my assessment is \textbf{moderately optimistic}, with substantial
uncertainty. I think an AI-assisted resolution in that interval is
plausible enough to justify a focused research effort. I would not treat
it as a likely deliverable on a fixed schedule. This is a subjective
assessment of the research situation, not a calibrated forecast of future
model capabilities. The evidence supports optimism about further useful
results more strongly than it supports a timetable for a full solution.

\paragraph{Reasons for optimism.}
This project already has an unusually useful starting point for future
models: a precise open question, a certified structural normal form, a
certified decidable fragment, explicit counterexamples to attractive
repairs, and a map of the current proof obligations. The fragment result
shows that a change in model construction can remove difficulties that
repeated local repairs had not resolved. A better model may find another
such change, recognize an applicable theorem in an adjacent field, or
discover a reduction whose required structure differs from the familiar
grid attempts. The short normalization and erasure arguments here also
show that there are still tractable structural consequences to extract
from the present framework; their modest scope should not be mistaken
for evidence that the final step is close.

Models with better mathematical reasoning, more reliable long-context
work, and better use of executable search and proof assistants could
explore these alternatives more systematically. In particular, finite
countermodel search can reject proposed local invariants cheaply, while
a proof assistant can prevent an accepted formal argument from silently
assuming its own completeness premise. A well-maintained record of failed
lemmas would help future attempts spend their effort on different ideas.

\paragraph{Reasons for caution.}
The remaining task is a finite-to-global existence or representation
problem, or an undecidability reduction. The project has not exposed a
missing routine calculation whose completion merely requires a larger
computation. A proposed finite state must still capture all relevant
interactions of an unbounded model. More fluent explanations, larger
bounded tests, and repeated agreement among similar models do not
establish that property. Closely related models can share the same
mistaken abstraction, so their agreement is weaker evidence than a
checked proof or a concrete counterexample.

Nor does the number of historical failures reliably measure how far the
problem is from a solution. A new invariant could shorten the argument
dramatically; equally, the surviving obstruction might resist several
further generations of models. The current evidence does not justify a
confident preference between eventual decidability and undecidability.
The decidable fragments give positive structure without resolving the
expressive power added by the unrestricted interactions.

\paragraph{What would increase my confidence.}
I would update the forecast substantially after a model produces one of
three concrete advances: a complete finite or automatic representation
theorem for the unrestricted semantics; a provably terminating procedure
whose soundness and completeness cover the full language; or a faithful
reduction from a known undecidable problem. The decisive intermediate
lemmas should be formalized early, with semantic adequacy and both
directions of any reduction reviewed independently. A polished proposed
proof that postpones its main completeness lemma would provide little
reason to update the forecast.

My practical expectation is therefore that increasingly capable AI,
combined with sustained human direction and formal checking, makes this
a worthwhile target for continued work. I am hopeful about a resolution
within a few years, but the present report supplies no mathematical basis
for promising one. The part of the workflow most likely to help is a
cycle of distinct structural hypotheses, adversarial counterexamples,
and early certification of the lemmas that survive.

\appendix
\section{Source map and finite checks}
\label{app:evidence}

The following references are to the supplied snapshot, not to an assumed
future state of the public repository. Line numbers are approximate
navigation aids; declaration names are the more stable identifiers.

\begin{center}
\small
\begin{tabularx}{\linewidth}{@{}p{0.47\linewidth}Y@{}}
\toprule
Current artifact or declaration & Location in the supplied source\\
\midrule
Ordered-disjoint normal form & Fragment paper, lines 318--443.\\
Cone definitions and correctness & Fragment paper, lines 444--693 and 814--917.\\
Full-language open premises & Overview, lines 2373--2418.\\
\code{coneScheme\_complete} & \code{POFreeLift.lean}, near line 42317.\\
\code{unf\_truth} & Same file, near lines 44006--44050.\\
\code{coneScheme\_sound} & Same file, near line 44093.\\
\code{decidableSat\_cone} & Same file, near line 44196.\\
\code{decidableFSat} & Same file, near line 44589.\\
\code{satisfiable\_iff\_set} & Same file, near line 44948.\\
\code{decidableSetSat} & Same file, near line 44976.\\
\bottomrule
\end{tabularx}
\end{center}

Two finite checks supported the review. First, enumerating all triples
of nonempty subsets of a four-element universe gave exactly the 25
composition-table cells printed in the current fragment paper. Second,
enumerating the 41 ordered-disjoint frames on three labelled points and
all eight valuations of $A$ tested the boundary formulas in
Section~\ref{sec:boundaries}. Counting satisfying combinations of frame,
valuation and designated root gave:

\begin{center}
\small
\begin{tabular}{@{}lrr@{}}
\toprule
Boundary concept & Erased formula & Original formula\\
\midrule
$C_{\exists\PO}$ & 480 & 0\\
$C_{\exists\PP}$ & 12 & 0\\
$C_{\forall\DR}$ & 12 & 0\\
\bottomrule
\end{tabular}
\end{center}

These counts are finite corroboration. The all-cardinality
unsatisfiability claims follow from the explicit relation case splits,
and the normalization and erasure theorems follow from their proofs,
not from bounded search. No kernel certification of the new results is
claimed.

\begin{thebibliography}{9}
\small
\bibitem{project}
\textit{ALCI\_RCC5 decidability project}.
Supplied snapshot \texttt{alcircc5-master(1).zip}, reviewed
7 October 2026. Current status sources: \path{README.md},
\path{approaches/ConeScheme/Overview.md}, and
\path{approaches/Lean_F6/Overview.md}.
Public repository: \url{https://github.com/lambdamikel/alcircc5}.

\bibitem{fragment}
\textit{Current focused paper on the $\forall\PO$-free fragment}.
In the supplied project snapshot:
\path{papers/pofree_fragment_arxiv.tex} and its PDF.

\bibitem{overview}
\textit{Current project overview paper}.
In the supplied project snapshot:
\path{papers/overview_arxiv.tex} and its PDF.

\bibitem{lean}
\textit{Current Lean development}.
In the supplied project snapshot: \path{formal/POFreeLift.lean}.
Full-language conditional pipeline:
\path{formal/Round19Transport.lean}.

\bibitem{lutzwolter}
C.~Lutz and F.~Wolter.
Modal logics of topological relations.
\textit{Logical Methods in Computer Science}, 2(2:5):1--41, 2006.
\url{https://doi.org/10.2168/LMCS-2(2:5)2006}.

\bibitem{wessel}
M.~Wessel.
\textit{Qualitative spatial reasoning with the ALCI\_RCC family:
first results and unanswered questions}.
Technical Report FBI-HH-M-324/03, Universit\"at Hamburg, 2002/2003.
Included in the supplied snapshot as \path{papers/report7.pdf}.
\end{thebibliography}
\end{document}
